\section{Rapport des tests du récepteur JSON}

\subsection{But}
Vérifier la validité d'un fichier JSON et l'intégrité des données reçues.

\subsection{Résumé des tests}
\begin{itemize}
  \item \hyperref[sec:json_example]{Étape 1 : créer un fichier JSON exemple}
  \item \hyperref[sec:json_receiver_code]{Étape 2 : créer un script de test qui importe le récepteur JSON et valide le JSON}
  \item \hyperref[sec:json_sender_code]{Étape 3 : créer un pseudo émetteur qui envoie le JSON au récepteur}
  \item \hyperref[sec:json_compare]{Étape 4 : comparer les données reçues avec les données envoyées}
\end{itemize}

\subsection{Tests}

\subsubsection{fichier json}
\label{sec:json_example}

Dans un premier temps, un fichier JSON exemple a été créé avec les données (spécifié lors de réunions) suivantes :
\begin{itemize}
    \item IMG : buffer de l'image en sortie de la caméra
    \item TEMPÉRATURE : température en degrés Celsius
    \item HUMIDITÉ : humidité en pourcentage
    \item LONGITUDE : longitude
    \item LATITUDE : latitude
    \item HORODATAGE : date et heure
    \item NIVEAU DE BATTERIE : niveau de batterie en pourcentage
\end{itemize}
or, d'autres champs ont été rajoutés pour les besoins de reconstruction d'image
les champs suivant ont été ajoutés :
\begin{itemize}
    \item TYPE : type de format de l'image (jpeg, png, etc.)
    \item WIDTH : largeur de l'image
    \item HEIGHT : hauteur de l'image
\end{itemize}

ce fichier json est disponible à : \texttt{include/code\_json\_receiver/data.json}

\subsubsection{Recepteur JSON}
\label{sec:json_receiver_code}

le récepteur JSON est disponible à : \texttt{include/code\_json\_receiver/receiver.py}

ce script permet de recevoir un objet JSON depuis une page web avec la méthode POST

\subsubsection{Émetteur de test}
\label{sec:json_sender_code}

le récepteur JSON est disponible à : \texttt{include/code\_json\_receiver/emitter.py}

ce script permet d'envoyer l'objet JSON (défini dans \hyperref[sec:json_example]{Étape 1}) à une page web avec la méthode POST

\subsubsection{Comparaison des données}
\label{sec:json_compare}

pour vérifier l'intégrité des données reçues, des JSON simples ont été créés.

JSON envoyé :
\begin{verbatim}
{
  "TEMP": 25,
  "HUMIDITY": 60
}
\end{verbatim}

JSON reçu :
\begin{verbatim}
{
  "TEMP": 25,
  "HUMIDITY": 60
}
\end{verbatim}

\subsection{Validation}

les données reçues sont identiques aux données envoyées, le test est donc concluant.

\subsection{Remarques}
\begin{itemize}
  \item Le récepteur JSON a correctement validé les fichiers JSON envoyés.
\end{itemize}



